% Jacquet，第3节后半；原书第71--75页。由23-jacquet-l-nonarch-factor.tex接入。
\origsubsection{3.3}{Langlands 构造（Langlands construction）}
下面回顾 Silberger 和 Wallach 将 Langlands 的结果推广到 $p$ 进情形的工作。

设 $Q=M_QU_Q$ 为型 $(p_1,p_2,\ldots,p_r)$ 的抛物子群（parabolic subgroup），并对每个 $1\leq i\leq r$，取 $G_{m_i}$ 的一个不可约本质缓增表示（irreducible essentially tempered representation）$\tau_i$。置 $\tau=\mathop{\times}_i\tau_i$。于是 $\tau_i=\tau_{i,0}\otimes\alpha^{t_i}$，其中 $\tau_{i,0}$ 缓增，$t_i$ 为实数。假定
\begin{equation}
t_1>t_2>\cdots>t_r.\tag{3.3.1}
\end{equation}
则诱导表示（induced representation）
\begin{equation}
\eta=I(G,Q;\tau_1,\tau_2,\ldots,\tau_r)\tag{3.3.2}
\end{equation}
有一个最大的真子表示 $\eta'$（它可以是 $\{0\}$）；不可约表示 $\pi=\eta/\eta'$ 记作
\begin{equation}
\pi=J(G,P;\tau_1,\tau_2,\ldots,\tau_r).\tag{3.3.3}
\end{equation}
每个不可约表示 $\pi$ 都有 (3.3.3) 的形式，其中（标准的）$P$ 及各 $\tau_i$ 由 $\pi$ 唯一确定。\footnote{译注：扫描页在本段先以 $p_i$ 表示分块大小，随后写作 $G_{m_i}$；式 (3.3.2) 用 $Q$，而式 (3.3.3) 及其后一句用 $P$。这里保留可见记号。}

$\eta'$ 的空间 $V'$ 可以显式描述。设 $W$ 为 $\tau$ 的空间，$\widetilde W$ 为逆反表示（contragredient representation）$\widetilde\tau$ 的空间。则 $V'$ 由 $\xi$ 的空间 $V$ 中满足下式的 $F$ 组成：
\begin{equation}
\int\langle F(\bar u g),\widetilde v\rangle\,d\bar u=0,
\qquad \widetilde v\in\widetilde W,\quad g\in G_n.\tag{3.3.4}
\end{equation}
此积分绝对收敛，积分域为 $\bar U=\bar U_Q=U_{\bar Q}$，其中 $\bar Q={}^tQ$ 是与 $Q$ 相反的抛物子群（opposite parabolic subgroup）。\footnote{译注：此处表示空间所属的表示在扫描页写作 $\xi$，前文相应诱导表示记为 $\eta$。}容易推出，$\pi$ 的矩阵系数（coefficient）可由与 (2.2.4) 类似的积分得到。具体地，设 $H:G_n\times G_n\longrightarrow\C$ 满足下列性质：
\begin{equation}
\begin{gathered}
H(u_1mg_1,\bar u_2mg_2)=H(g_1,g_2),\\
u_1\in U_Q,\quad\bar u_2\in\bar U_Q,\quad m\in M_Q;
\end{gathered}\tag{3.3.5}
\end{equation}
\begin{equation}
\begin{gathered}
\text{对任意 }g_1,g_2,\quad m\longmapsto H(mg_1,g_2)\\
\text{是 }\tau\otimes\delta_Q^{1/2}\text{ 的矩阵系数；}
\end{gathered}\tag{3.3.6}
\end{equation}
\begin{equation}
H\text{ 关于右作用是 }K_n\times K_n\text{ 有限的。}\tag{3.3.7}
\end{equation}
则由收敛积分
\begin{equation}
f(g)=\int_{M\backslash G}H(hg,h)\,dh
=\int_{\bar U\times K}H(\bar u kg,k)\,dk\,d\bar u\tag{3.3.8}
\end{equation}
定义的函数 $f$ 是 $\pi$ 的矩阵系数，并且 $\pi$ 的所有矩阵系数都可通过适当的 $H$ 以此方式得到（关于收敛问题，见 (3.6.6) 和 (3.6.7)）。

可以用 $\bar Q$ 代替 $Q$。此时，只要将条件 (3.3.1) 换成不等号方向相反的类似条件，就能定义与 (3.3.3) 相应的表示。例如，商表示
\begin{equation}
\pi'=J(G,\bar Q;\widetilde\tau_1,\widetilde\tau_2,\ldots,\widetilde\tau_r)\tag{3.3.9}
\end{equation}
是有定义的，它是诱导表示
\begin{equation}
\eta'=I(G,\bar Q;\widetilde\tau_1,\widetilde\tau_2,\ldots,\widetilde\tau_r)\tag{3.3.10}
\end{equation}
的商。其矩阵系数由积分 (3.3.8) 给出，其中 $H$ 满足 (3.1.5)--(3.1.7)，并将 $(Q,\tau)$ 换成 $(\bar Q,\widetilde\tau)$。\footnote{译注：扫描页此处交叉引用可见为 (3.1.5)--(3.1.7)；上文对应的三个条件编号为 (3.3.5)--(3.3.7)。}

特别地，设 $f$ 为由 (3.3.8) 定义的 $\pi$ 的矩阵系数，其中 $H$ 对 $(Q,\tau)$ 满足 (3.3.5)--(3.3.7)。则
\begin{equation}
f^\vee(g)=\int_{M\backslash G}\widetilde H(hg,h)\,dh,\tag{3.3.11}
\end{equation}
其中 $\widetilde H(g_1,g_2)=H(g_2,g_1)$。而 $\widetilde H$ 确实对 $(\bar Q,\widetilde\tau)$ 满足 (3.3.5)--(3.3.7)。所以 $f^\vee$ 是 $\pi'$ 的矩阵系数，且
\begin{equation}
\pi'=\widetilde\pi.\tag{3.3.12}
\end{equation}
当然，$\bar Q$ 也与型 $(n_r,n_{r-1},\ldots,n_1)$ 的标准抛物子群 $Q'$ 共轭，因而
\begin{equation}
\widetilde\pi=\pi'=J(G,Q';\widetilde\tau_r,\widetilde\tau_{r-1},\ldots,\widetilde\tau_1).\tag{3.3.13}
\end{equation}

\begin{originaltheorem}{定理 3.4（Theorem 3.4）}
沿用 (3.3) 的记号。则
\[
L(s,\pi)=\prod_iL(s,\tau_i),\qquad
L(s,\widetilde\pi)=\prod_iL(s,\widetilde\tau_i),
\]
\[
\varepsilon(s,\pi,\psi)=\prod_i\varepsilon(s,\tau_i,\psi).
\]
\end{originaltheorem}
\begin{proof}
只需证明关于 $L(s,\pi)$ 的断言。事实上，交换 $\pi$ 与 $\widetilde\pi$ 即可得到关于 $L(s,\widetilde\pi)$ 的断言；最后一个断言则由 (2.7) 和关系
\[
\gamma(s,\eta,\psi)=\prod_i\gamma(s,\tau_i,\psi)
\]
（参见 (2.3)）推出。证明分为本小节余下的论证和 (3.5) 中的构造两部分。由 (2.3) 和 (2.7)，存在多项式 $P$ 和 $\widetilde P$，使
\[
L(s,\pi)=P(q^{-s})\prod_iL(s,\tau_i),\qquad
L(s,\widetilde\pi)=\widetilde P(q^{-s})\prod_iL(s,\widetilde\tau_i).
\]
由 (3.2.4) 与 (3.3.1) 可知，$L(s,\tau_1)^{-1}$ 与 $\prod_iL(1-s,\widetilde\tau_i)^{-1}$ 互素。因此，$P(q^{-s})$ 与 $L(s,\tau_1)^{-1}$ 互素（(2.7.5)）。所以，若 $P\ne1$，则
\begin{equation}
\prod_{2\leq i\leq r}L(s,\tau_i)\tag{3.4.1}
\end{equation}
不属于 $I(\pi)$。于是只需证明 (3.4.1) 实际上属于 $I(\pi)$。

\paragraph*{3.5.}
恒等式
\begin{equation}
L(s,\pi)=\prod_{1\leq i\leq r}L(s,\tau_i)\tag{3.5.1}
\end{equation}
在 $r=1$ 时显然成立。假定它对 $r-1$ 成立，来证明它对 $r$ 成立。如刚才所见，只需证明 (3.4.1) 属于 $I(\pi)$。

置 $n_1=p_1$，$n_2=p_2+\cdots+p_r$。设 $Q'$ 为 $G_{n_2}$ 中型 $(p_2,\ldots,p_r)$ 的抛物子群。于是表示
\begin{equation}
\sigma_2=J(G_{n_2},Q';\tau_2,\ldots,\tau_r)\tag{3.5.2}
\end{equation}
是有定义的。再置 $\sigma_1=\tau_1$、$\sigma=\sigma_1\times\sigma_2$，设 $P$ 为 $G_n$ 中型 $(n_1,n_2)$ 的抛物子群。则 $\pi$ 实际上是
\begin{equation}
\xi=I(G,P;\sigma_1,\sigma_2)\tag{3.5.3}
\end{equation}
的一个商表示。更精确地，容易看出，$\pi$ 的矩阵系数由积分 (3.3.8) 给出，其中 $H$ 为对 $(P,\sigma)$ 满足 (3.3.5)--(3.3.7) 的任意函数，即用 $(P,\sigma)$ 代替 $(Q,\tau)$。

若 $f$ 是以这种方式定义的 $\pi$ 的矩阵系数，则
\begin{equation}
\begin{aligned}
&Z\bigl(\Phi,s+\tfrac12(n-1),f\bigr)\\
&\quad=\iint dk\,d\bar u\int_{G_n}|\det g|^{s+(n-1)/2}
\Phi(k^{-1}\bar u^{-1}g)H(g,k)\,d^\times g,
\end{aligned}\tag{3.5.4}
\end{equation}
其中 $\bar u$ 的积分域为 $\bar U=U_{\bar P}$。这也等于
\begin{equation}
\begin{aligned}
&\iint dk\,dk'\iint
|\det m_1|^{s+(n_1-1)/2}|\det m_2|^{s+(n_2-1)/2}\\
&\quad\cdot\delta_{\bar P}\!\begin{pmatrix}m_1&0\\0&m_2\end{pmatrix}^{-1/2}
H\!\left[\begin{pmatrix}m_1&0\\0&m_2\end{pmatrix}k',k\right]
\,d^\times m_1\,d^\times m_2\\
&\quad\cdot\iint\Phi\!\left[k^{-1}\begin{pmatrix}m_1&y\\xm_1&m_2+xy\end{pmatrix}k'\right]\,dx\,dy.
\end{aligned}\tag{3.5.5}
\end{equation}
这里 $m_i$ 遍历 $G_{n_i}$。我们将看到，给定 $\Phi_2\in\mathcal S(n_2\times n_2,F)$ 和 $\sigma_2$ 的矩阵系数 $f_2$，存在 $H$ 和 $\Phi$，使 $Z(\Phi_2,s+\tfrac12(n_2-1),f_2)$ 等于 (3.5.5)。由于根据归纳假设，$L(s,\sigma_2)$ 等于 (3.4.1)，由此便有 (3.4.1) 属于 $I(\pi)$，证明也就完成了。

存在 $\pi_1$ 的一个矩阵系数 $f_1$ 和一个 $\Phi_1\in\mathcal S(n_1\times n_1,F)$，后者的支集包含在 $e$ 的某个紧邻域内，使得
\begin{equation}
Z\bigl(\Phi_1,s+\tfrac12(n_1-1),f_1\bigr)=1.\tag{3.5.6}
\end{equation}
还存在 $\Phi_{12}\in\mathcal S(n_1\times n_2,F)$ 和 $\Phi_{21}\in\mathcal S(n_2\times n_1,F)$，它们的支集包含在 $0$ 的某个邻域内，使得由
\begin{equation}
\Phi\!\begin{pmatrix}m_1&y\\x&m_2\end{pmatrix}
=\Phi_1(m_1)\Phi_{12}(y)\Phi_{21}(x)\Phi_2(m_2)\tag{3.5.7}
\end{equation}
定义的 $\Phi\in\mathcal S(n\times n,F)$ 满足
\begin{equation}
\iint\Phi\!\begin{pmatrix}m_1&y\\xm_1&m_2+xy\end{pmatrix}\,dx\,dy
=\Phi_1(m_1)\Phi_2(m_2).\tag{3.5.8}
\end{equation}
还存在 $K$ 上两个 $K$ 有限函数（$K$-finite functions）$\xi$ 和 $\xi'$，使得
\begin{equation}
\iint\Phi(k^{-1}zk')\xi(k)\xi'(k')\,dk\,dk'=\Phi(z).\tag{3.5.9}
\end{equation}
于是
\begin{equation}
\begin{aligned}
&Z\bigl(\Phi_2,s+\tfrac12(n-1),f_2\bigr)\\
&\quad=\iint|\det m_1|^{s+(n_1-1)/2}|\det m_2|^{s+(n_2-1)/2}
 f_1(m_1)f_2(m_2)\,d^\times m_1\,d^\times m_2\\
&\qquad\cdot\iint\xi(k)\xi'(k')\,dk\,dk'
\iint\Phi\!\left[k^{-1}\begin{pmatrix}m_1&y\\xm_1&m_2+xy\end{pmatrix}k'\right]\,dx\,dy.
\end{aligned}\tag{3.5.10}
\end{equation}


最后沿用 (2.5) 的平均化方法完成证明。\footnote{译注：扫描页在引入 $f_1$ 时写 $\pi_1$，本段上文相应表示为 $\sigma_1$；式 (3.5.10) 左端写 $s+\tfrac12(n-1)$，而前文同一 $\Phi_2,f_2$ 的积分写作 $s+\tfrac12(n_2-1)$。式 (3.5.5) 的模特征下标可见为 $\bar P$，而式 (3.5.11) 为 $P$。均按可见原式保留。}具体地，设 $dh$ 和 $dh'$ 分别为 $K\cap\bar P$ 和 $K\cap P$ 上的归一化 Haar 测度（normalized Haar measure）。对 $h\in K\cap\bar P$、$h'\in K\cap P$，写
\[
h=\begin{pmatrix}h_1&0\\ *&h_2\end{pmatrix},\qquad
h'=\begin{pmatrix}h'_1&*\\0&h'_2\end{pmatrix}.
\]
则 (3.5.10) 还等于
\begin{equation}
\begin{aligned}
&\iint dk\,dk'\iint|\det m_1|^{s+(n_1-1)/2}|\det m_2|^{s+(n_2-1)/2}\\
&\quad\cdot\delta_P\!\begin{pmatrix}m_1&0\\0&m_2\end{pmatrix}^{-1/2}
H_1\!\left[\begin{pmatrix}m_1&0\\0&m_2\end{pmatrix},k,k'\right]
\,d^\times m_1\,d^\times m_2\\
&\quad\cdot\iint\Phi\!\left[k^{-1}\begin{pmatrix}m_1&y\\xm_1&m_2+xy\end{pmatrix}k'\right]\,dx\,dy,
\end{aligned}\tag{3.5.11}
\end{equation}
其中
\begin{equation}
\begin{aligned}
&H_1\!\left[\begin{pmatrix}m_1&0\\0&m_2\end{pmatrix},k,k'\right]\\
&\quad=\delta_P\!\begin{pmatrix}m_1&0\\0&m_2\end{pmatrix}^{1/2}
\iint\xi(hk)\xi'(h'k')f_1(h_1m_1h_1'^{-1})f_2(h_2m_2h_2'^{-1})\,dh\,dh'.
\end{aligned}\tag{3.5.12}
\end{equation}
存在对 $(P,\sigma)$ 满足 (3.3.5)--(3.3.7) 的 $H$，使
\begin{equation}
H\!\left[\begin{pmatrix}m_1&0\\0&m_2\end{pmatrix}k',k\right]
=H_1\!\left[\begin{pmatrix}m_1&0\\0&m_2\end{pmatrix},k,k'\right].\tag{3.5.13}
\end{equation}
比较 (3.5.11) 与 (3.5.5)，便得到所要的断言。
\end{proof}

\origsubsection{3.6}{非分歧表示（Unramified representations）}
设 $\pi$ 非分歧，即含有非零的 $K$-不变向量（也即包含 $K$ 的平凡表示，unit representation）。用 $B_n$ 表示型 $(1,1,1,\ldots,1)$ 的抛物子群，则存在 $u_i\in\C$，使 $\pi$ 为下列诱导表示的唯一非分歧成分：
\begin{equation}
\xi=I(G,B_n;\mu_1,\mu_2,\ldots,\mu_n),\qquad \mu_i=\alpha^{u_i}.\tag{3.6.1}
\end{equation}
甚至可以假定 $u_1\geq u_2\geq\cdots\geq u_n$。改变记号，可将由各 $\mu_i$ 组成的 $n$ 元组写成
\[
(\nu_1^1,\nu_2^1,\ldots,\nu_{n_1}^1,
 \nu_1^2,\nu_2^2,\ldots,\nu_{n_2}^2,\ldots,
 \nu_1^r,\nu_2^r,\ldots,\nu_{n_r}^r),
\]
其中 $n=n_1+n_2+\cdots+n_r$，$\nu_i^j=\alpha^{t_j}\lambda_i^j$，$\lambda_i^j$ 是特征（character），$t_j$ 为实数，且 $t_1>t_2>\cdots>t_r$。于是
\[
\tau_j=I(G_{n_j},B_{n_j};\nu_1^j,\nu_2^j,\ldots,\nu_{n_j}^j)
\]
是 $G_{n_j}$ 的本质缓增非分歧表示；事实上，若 $Q$ 为型 $(n_1,n_2,\ldots,n_r)$ 的抛物子群，则
\begin{equation}
\pi=J(G,Q;\tau_1,\tau_2,\ldots,\tau_r).\tag{3.6.2}
\end{equation}
这由一个显式计算得出：若 $H$ 对 $(Q,\mathop{\times}_j\tau_j)$ 满足条件 (3.3.5)--(3.3.7)，且 $H(k,k')=1$，则
\begin{equation}
\int H(\bar u,e)\,d\bar u\ne0.\tag{3.6.3}
\end{equation}
所以，$\eta=I(G,Q;\tau_1,\tau_2,\ldots,\tau_r)$ 的“$J$ 成分”是非分歧的；由于 $\xi\simeq\eta$，便得到 (3.6.2)。\footnote{译注：扫描页先取 $u_i\in\C$，随后直接写 $u_1\geq\cdots\geq u_n$，未在该行另写实部；此处保留其可见形式。}

因此
\begin{equation}
\begin{gathered}
L(s,\pi)=\prod_iL(s,\mu_i),\qquad
L(s,\widetilde\pi)=\prod_iL(s,\mu_i^{-1}),\\
\varepsilon(s,\pi,\psi)=\prod_i\varepsilon(s,\mu_i,\psi)
\quad\text{（若 $\psi$ 的指数为零，则等于 $1$）。}
\end{gathered}\tag{3.6.4}
\end{equation}
更精确的结果见 [R.G.-H.J., \S6]。

\begin{originaltheorem}{备注 3.6.5（Remark 3.6.5）}
沿用 (2.3) 的记号。设各 $\sigma_i$ 不可约且非分歧。则 $\xi$ 有唯一的非分歧成分 $\pi$。由 (3.6.4) 容易推出
\[
L(s,\pi)=\prod_iL(s,\sigma_i),
\]
而 $L(s,\widetilde\pi)$ 和 $\varepsilon$ 也有类似的关系。
\end{originaltheorem}

\begin{originaltheorem}{备注 3.6.6（Remark 3.6.6）}
设 $\pi$ 是 $G_n$ 的缓增表示。令 $\rho_n=I(G_n,B_n;1,1,\ldots,1)$，以 $\Xi_n$ 表示附属于 $\rho$ 的球函数（spherical function）。则 $\pi$ 的每个矩阵系数的绝对值都以 $\Xi_n$ 的某个常数倍为上界。再沿用 (3.3) 的记号，则 $\tau_i$ 的矩阵系数的绝对值都以 $\Xi_{n_i}\otimes\alpha^{t_i}$ 的某个常数倍为上界。此外，若 $H$ 对 $(Q,\tau)$ 满足 (3.3.5)--(3.3.7)，则 $|H|\leq cH_0$，其中 $H_0$ 满足 (3.3.5)，且
\[
H_0(mk,k')=\delta_P^{1/2}(m)\prod_i\Xi_{n_i}(m_i).
\]
所以，就收敛问题而言，可以将 $\tau_i$ 换成 $\rho_{n_i}\otimes\alpha^{t_i}$，并将 $H$ 换成 $H_0>0$。结合下一个引理，就解决了所有收敛问题。\footnote{译注：本显示式在扫描页的下标可见为 $P$，且乘积内未另写前文的 $\alpha^{t_i}$ 因子；这里保留原式。}
\end{originaltheorem}

\begin{originaltheorem}{引理 3.6.7（Lemma 3.6.7）}
若 $\Omega$ 是 $\mathcal S(n\times n,F)$ 的一个有界集（bounded set），则存在 $\mathcal S(n\times n,F)$ 中的 $\Phi_0\geq0$，使得对 $k_i\in K$ 有 $\Phi_0(k_1xk_2)=\Phi_0(x)$，且对 $\Phi\in\Omega$ 有 $|\Phi|\leq\Phi_0$。
\end{originaltheorem}
\begin{proof}
可以假定，$\Omega$ 还包含所有由 $\Phi\in\Omega$ 得到的函数 $x\mapsto\Phi(k_1xk_2)$。于是存在 $\Phi_1\geq0$，使得对 $\Omega$ 中所有的 $\Phi$ 都有 $|\Phi|\leq\Phi_1$ [A.W. 1, Lemma 5]。只需取
\[
\Phi_0(x)=\iint\Phi_1(k_1xk_2)\,dk_1\,dk_2,
\]
其中 $dk_i$ 是 $K$ 上的归一化 Haar 测度。
\end{proof}

\paragraph*{3.7.}
按照“Langlands 哲学（Langlands' philosophy）”，Weil--Deligne 群（Weil--Deligne group）$W'_F$ 的 $n$ 维半单表示（semisimple representation）与 $G_n(F)$ 的不可约容许表示（irreducible admissible representation）之间，应当存在一个“自然双射（natural bijection）”$\sigma\mapsto\pi(\sigma)$。此外，对 $\pi=\pi(\sigma)$，还应有
\[
L(s,\pi)=L(s,\sigma),\qquad
\varepsilon(s,\pi,\psi)=\varepsilon(s,\sigma,\psi),
\]
\[
\pi(\sigma)^\sim=\pi(\widetilde\sigma),\qquad
\pi(\sigma\otimes\chi)=\pi(\sigma)\otimes\chi.
\]
显然，若能对 Weil 群（Weil group）的不可约表示定义映射 $\sigma\mapsto\pi(\sigma)$，则这一猜想就能得到证明（参见下文第5节）。
